\begin{claim}{C.58}{Proved}{A uniform word bound for the full product affine group}
Let
\[
 A=\operatorname{Aff}^+(\R)\times\operatorname{Aff}^+(\R).
\]
For a continuous path $k:[0,t]\to A$ with $k(0)=1$, put
$P=k([0,t])\cup k([0,t])^{-1}$ and $g=k(t)$.
Every exponential fraction satisfies the exact inclusion
\[
 \{\exp(r\log g):0\leq r\leq1\}\subseteq P^M,
 \qquad M=172\,496\,624\,025\,607.                       \tag{C.58.1}
\]
The integer is independent of the path and all scale coordinates,
including zero coordinates.

In the plane-loop setting of C.54--C.56, let $U^\flat$
be the union of $U^\sharp$ and all the remaining C.54
open orbits, whose actual generated source is $A$, and set
$X^\flat=\overline{U^\flat}$.
The prescribed source exponential fractions extend to a
jointly continuous based contraction through homeomorphisms
of $X^\flat$, with jointly continuous inverse and exact
$H$-commutation. They agree with the contraction on
$X^\sharp$ and retain the based exponential prescription
on every nonopen orbit in $X^\flat$.
The finite-representative extension of C.55 also holds
relative to $X^\flat$, for representatives outside that
set having compact quotients by their full local centralizers.
\end{claim}

\begin{proof}
\emph{Seven signed values in logarithmic coordinates.}
For every continuous based path $\gamma:[0,t]\to\R^4$
and $0\leq r\leq1$, the signed-measure result
\cite[Theorem~4.7 and \S4.3.2]{C_GrandoniRaviSinghZenklusen2014}
gives
\[
 r\gamma(t)=\sum_{j=1}^{m}\varepsilon_j\gamma(s_j),
 \quad m\leq7,\quad \varepsilon_j\in\{1,-1\},\quad
                         0\leq s_j\leq t.              \tag{C.58.2}
\]
For a polygonal path, apply it to the four coordinate
derivatives. Three circular intervals contribute six
signed endpoint values, and at most one wrap contributes
the total endpoint. The required halving count is explicit:
three intervals cut at four points give at most seven
pieces, so one half uses at most three intervals.
For the whole interval four cuts give at most five pieces,
so one half uses at most two. Complementation, halving
and compactness supply all fractions.
Uniform polygonal approximation and compactness of seven
times and their signs give (C.58.2) for continuous paths.

Apply (C.58.2) to the actual path $\log k(s)$, using the
global logarithm of C.54. Pad by $\log k(0)=0$ so that
$m=7$, and put
\[
 X_j=\varepsilon_j\log k(s_j),\qquad
 \ell_j=\exp X_j=k(s_j)^{\varepsilon_j},\qquad
 w=\ell_1\cdots\ell_7.
                                                               \tag{C.58.3}
\]
Thus $\sum_jX_j=r\log g$. All inverse letters are included
in the signed Lie coordinates $X_j$; the equality
$\exp(-X)=(\exp X)^{-1}$ is the one-parameter group identity.

\emph{The full product error and its antisymmetric identity.}
In one affine factor use logarithmic scale and group
translation coordinates
\[
 (a,v)(b,z)=(a+b,v+e^az),\qquad
 \sigma(a)=\int_0^1e^{ua}\,du,\qquad
                  \exp(a,u)=(a,\sigma(a)u).             \tag{C.58.4}
\]
The function $\sigma$ is positive and smooth, with
$\sigma(0)=1$. Write $X_j=(a_j,u_j)$ in this factor,
and set $T=\sum_j a_j$ and $P_j=\sum_{\ell<j}a_\ell$.
The pure translation taking $w$ to $\exp(\sum_jX_j)$
by \emph{left} multiplication is
\[
 c=\sum_{j=1}^7 F_j(a)u_j,\qquad
 F_j(a)=\sigma(T)-e^{P_j}\sigma(a_j),\qquad
                         a=(a_1,\ldots,a_7).           \tag{C.58.5}
\]
Indeed the product translation is
$\sum_je^{P_j}\sigma(a_j)u_j$, whereas the target
translation is $\sigma(T)\sum_ju_j$.
The exponential telescoping identity proves
\[
 \sum_j a_jF_j(a)
 =T\sigma(T)-\sum_je^{P_j}(e^{a_j}-1)=0.               \tag{C.58.6}
\]

For all $i,j$ define the smooth antisymmetric functions
\[
 f_{ij}(a)=\int_0^1 s\,
       (\partial_iF_j-\partial_jF_i)(sa)\,ds.
                                                               \tag{C.58.7}
\]
Differentiating (C.58.6) at $sa$ gives
$\sum_i a_i\partial_jF_i(sa)=-F_j(sa)/s$ for $s>0$.
Consequently
\[
 s\sum_i a_i(\partial_iF_j-\partial_jF_i)(sa)
                   =\frac{d}{ds}\{sF_j(sa)\}.
 \]
The endpoint at zero vanishes by continuity.
Integration therefore gives $F_j=\sum_i a_if_{ij}$ and
\[
 c=\sum_{i<j}f_{ij}(a)(a_i u_j-a_j u_i).                \tag{C.58.8}
\]
An explicit form, including its continuous interpretation
at $a_j=0$, is
\[
 \begin{split}
 f_{ij}(a)
   &=-\int_0^1s e^{sP_j}\sigma(sa_j)\,ds\\
   &=-\frac{\sigma(P_j+a_j)-\sigma(P_j)}{a_j}
          \quad(i<j),\qquad
 f_{ij}\big|_{a_j=0}=-\sigma'(P_j).
 \end{split}                                            \tag{C.58.9}
\]
Here the two derivatives of the common term $\sigma(T)$
cancel, and only $\partial_i(e^{P_j})$ remains.
Formula (C.58.9) also verifies (C.58.8) by telescoping.

With $[x,y]=xyx^{-1}y^{-1}$, direct multiplication gives
\[
 [\ell_i,\ell_j]
  =\bigl(0,\sigma(a_i)\sigma(a_j)
                         (a_i u_j-a_j u_i)\bigr).
                                                               \tag{C.58.10}
\]
Thus, putting
$G_{ij}(a)=f_{ij}(a)/(\sigma(a_i)\sigma(a_j))$,
the exact error in (C.58.5) is a sum of actual
commutator translations multiplied by $G_{ij}$.
In the product group the two factors evaluate
\emph{this same function} at their two scale vectors
$a^{(1)},a^{(2)}\in\R^7$.

\emph{Bounds smooth across every zero scale.}
Let $R=\sum_\ell|a_\ell|$.
The integral expression for $\sigma$ bounds it and its
first two derivatives at $x$ by $e^{|x|}$.
Each $F_j$ and each of its first and second partial
derivatives therefore have absolute value at most $2e^R$.
Equation (C.58.7) yields
\[
 |f_{ij}(a)|\leq2e^R,\qquad
 |\partial_\ell f_{ij}(a)|\leq\tfrac43e^R.
 \]
Also
$\sigma(x)\geq e^{-|x|}$ and
$0\leq\sigma'(x)/\sigma(x)\leq1$.
Since $i\ne j$, differentiation in one coordinate
affects at most one of the two denominator factors.
It follows that
\[
 |G_{ij}(a)|\leq2e^{2R},\qquad
                 |\partial_\ell G_{ij}(a)|<4e^{2R}.
                                                               \tag{C.58.11}
\]
No nonzero scale assumption occurs in these estimates.

\emph{Bounded interpolation at the two seven-dimensional nodes.}
For any one of the functions $G=G_{ij}$ put
\[
 \begin{split}
 \mathcal A(a)&=\prod_{\ell=1}^7
                         (e^{a_\ell}+e^{-a_\ell})^4,\\
 \mathcal B_q(a)&=\mathcal A(a)\tanh a_q
                                  \quad(1\leq q\leq7),\\
 z_\ell&=\tanh a_\ell,\qquad
                         \phi(z)=G(a)/\mathcal A(a).
 \end{split}                                            \tag{C.58.12}
\]
The map $a\mapsto z$ is a diffeomorphism onto $(-1,1)^7$.
We have $\mathcal A\geq e^{4R}$ and
$\partial_\ell\log\mathcal A=4\tanh a_\ell$.
Equation (C.58.11) gives
\[
 |\phi|\leq2,\quad
 |\partial_{a_\ell}(G/\mathcal A)|\leq12e^{-2R},
 \quad |\partial_{z_\ell}\phi|\leq12.                  \tag{C.58.13}
\]
For the last inequality use
$\operatorname{sech}^2a_\ell\geq e^{-2|a_\ell|}$.

For unequal nodes $z^{(1)},z^{(2)}$, choose $q$ maximizing
$|z_q^{(1)}-z_q^{(2)}|$, and define
\[
 d=\frac{\phi(z^{(1)})-\phi(z^{(2)})}
           {z_q^{(1)}-z_q^{(2)}},\qquad
                         b=\phi(z^{(1)})-d z_q^{(1)}.
 \]
Integration along the segment between the nodes in the
convex cube $(-1,1)^7$ gives $|d|\leq12\cdot7=84$
and $|b|\leq86$.
For equal nodes take $d=0$ and $b=\phi(z^{(1)})$,
with any $q$. In both cases the exact two equations are
\[
 G(a^{(\nu)})=
     b\mathcal A(a^{(\nu)})+d\mathcal B_q(a^{(\nu)})
 \quad(\nu=1,2),\qquad |b|\leq86,\quad |d|\leq84.
                                                               \tag{C.58.14}
\]
In particular, independent affine scales cause no
unbounded divided difference: the complete scale vectors
are the nodes of the same smooth function.

\emph{Realization by actual translation paths.}
For a fixed pair $i<j$, the path
\[
 C_{ij}(v)=
 [k(vs_i)^{\varepsilon_i},k(vs_j)^{\varepsilon_j}],
                    \qquad 0\leq v\leq1,               \tag{C.58.15}
\]
starts at the identity and lies in the translation
plane of $A$. Its values and inverses use at most
four letters of $P$.
Put $Z_\ell=e^{a_\ell}$ in (C.58.12).
Both $\mathcal A$ and $\mathcal B_q$ are integer Laurent
polynomials in seven variables.
Their exponents have absolute value at most four in
each variable, and both sums of absolute coefficients
are at most $2^{28}$.
Indeed each factor of $\mathcal A$ has coefficient sum
$2^4$, while the distinguished factor of $\mathcal B_q$
is $(Z_q-Z_q^{-1})(Z_q+Z_q^{-1})^3$ with coefficient
norm at most $2^4$.

For a Laurent monomial $Z_1^{n_1}\cdots Z_7^{n_7}$ use
the conjugator $h_n=\ell_1^{n_1}\cdots\ell_7^{n_7}$.
It uses at most $28$ letters.
Conjugation by $h_n$ multiplies a translation's
$\nu$th coordinate by
$\exp(\sum_\ell n_\ell a_\ell^{(\nu)})$, irrespective
of the conjugator's own translation.
Expand either polynomial and multiply the corresponding
integer numbers of conjugated copies of $C_{ij}(v)$
or its inverse.
The resulting continuous based translation paths
$\mathcal A_{ij}(v)$ and $\mathcal B_{ij,q}(v)$
multiply the $\nu$th coordinate of $C_{ij}(v)$ by
$\mathcal A(a^{(\nu)})$ and
$\mathcal B_q(a^{(\nu)})$, respectively.
All their factors are translations and commute.
Every value or inverse has the uniform cost
\[
                     L=(2\cdot28+4)2^{28}=60\cdot2^{28}.
                                                               \tag{C.58.16}
\]

The planar signed-path identity of C.37 represents a
fraction of either translation endpoint by three signed
values of its actual path.
Use $86$ copies of the fraction $|b|/86$, with inverse
for negative $b$, and $84$ copies for $|d|/84$.
The resulting translation
$\mathcal A_{ij}(1)^b\mathcal B_{ij,q}(1)^d$ has cost
at most $3(86+84)L$ and, by (C.58.14), has exactly the
two coordinate corrections
$G_{ij}(a^{(\nu)})\,\operatorname{pr}_{v_\nu}C_{ij}(1)$.
There are at most $\binom72=21$ pairs.
Their translations commute, and their product followed
by $w$ is exactly $\exp(\sum_jX_j)=g^r$.
Thus the total cost is
\[
 7+21\cdot3(86+84)\cdot60\cdot2^{28}
                         =172\,496\,624\,025\,607.      \tag{C.58.17}
\]
Zero path values pad shorter words.
The cuts and interpolation choices need not vary
continuously; all paths used to obtain their words are
the actual continuous paths just specified.
The equality is in the original group $A$, with the
full product error corrected on the left.

\emph{Exact executable checks.}
The accompanying program
\path{code/part_b/c58_affine_product_word_certificate.py}
verifies the product, inverse and commutator formulas,
the explicit radial coefficient and its telescoping
identity, the zero-scale values, the two-node equations,
the Laurent coefficient norms and the integer bound.
It also checks the entire corrected product for every
signed pattern of lengths one through seven in several
two-factor scale regimes, retaining zero and canceling
prefix scales by exact formal Laurent arithmetic.
The signed-measure selection, analytic estimates and
topological conclusions have the proofs in this text;
the finite checks verify their computable algebra.

\emph{Arbitrary accumulation and the remaining relative statement.}
C.54 supplies the actual closed simply connected source
group on each new orbit, its entire closure action, its
operator transport and every global quotient equation.
C.56 supplies the old contraction on $X^\sharp$.
The bound in (C.58.17) exceeds all the earlier bounds.
Therefore all source fractions on $U^\flat$ belong to
one compact family of actual global word isotopies,
obtained by replacing the length in C.50.15 by $M$.
This is a finite union of continuous images of the
compact cube of word parameters; maps and inverses
are both included. Each endpoint equality holds in
the specified simply connected source, including the
based lifts in the older two-dimensional cases.

All these sources are metabelian. Their group identity
passes from $U^\flat$ to $X^\flat$, so the nontrivial
nonopen intrinsic images there are precisely the line
translation, line affine and circle rotation images
of C.37. On them use the unique line logarithm or the
specified based real circle angle, and fix point orbits.

For a converging sequence of source fractions
$w_i=\exp(r_i\log g_i)$, every positive integer $n$
gives the exact source residual
\[
 w_i^n g_i^{-\lfloor nr_i\rfloor}
       =\exp\bigl((nr_i-\lfloor nr_i\rfloor)\log g_i\bigr).
                                                               \tag{C.58.18}
\]
The residual fraction lies in $[0,1)$ and has the
same fixed bound $M$.
The compact word-parameter argument of C.50 now applies
without a change of topology or a new boundary prescription.
On a line the limit of the residual logarithms is bounded,
and the limiting $w$ commutes with $g$, by passing the
source identity to every point of the target orbit.
Dividing the logarithmic residual equation by $n$
forces $\log w=r\log g$.

For completeness, the circle comparison retains the
real lift. If $A_W,T,C_n$ are the real endpoint angles
of the limiting word, original path and residual word,
respectively, the equality is
$nA_W-m_nT=C_n$, with $|m_n-nr|\leq1$ and
$|C_n|\leq M\max_u|\theta(u)|$.
The source loop
$W_i(u)^n f_{ut_i}^{-m_i}Z_i(u)^{-1}$
is null in its specified simply connected source,
by (C.58.18). Evaluate it at two points converging to
distinct points of the target circle.
Its difference-direction winding is zero and stays
zero in the limit. C.29 identifies that winding
with the integer ambiguity in the displayed real
angle equation. Thus $A_W=rT$, including an identity
circle endpoint with nonzero specified real angle.
At a fixed point every limiting word is the identity.

These are exactly the dense-set limits required by
C.50's diagonal approximation on the closure.
They prove joint continuity on $X^\flat$.
The inverse loop has the same sources and quotient
families, and its fractions are the orbitwise inverses,
so it supplies a jointly continuous inverse family.
The based-edge equations and operator commutation pass
from the dense union to its closure.
Agreement on $U^\sharp$ gives agreement on $X^\sharp$.

Finally, an outside representative in the last
assertion is also outside $X^\ddagger$.
C.55's direct equivariant-map evaluation and its
compact fundamental set therefore give the same
whole-closure limit for its intrinsic contraction.
Its intersection with $X^\flat$ consists of nonopen
orbits. The actual group image on that intersection
is metabelian, so its hybrid target value is exactly
the exponential prescription just proved.
Finite closed-set pasting of both families and their
inverses gives the asserted relative extension.

This includes all C.54 source groups through arbitrary
accumulation. It does not establish the remaining
relative extension through the noncompact local
quotients outside this region, through infinitely
many unrelated residual representatives, or through
the exterior nonopen locus. C.12 remains open.
\end{proof}
